\section{Gluing kernels and stability from a fixed core}\label{sec:kernels}

We now return to the bare height weights, with no confining penalty.
The flat-pin estimate gives an average bound for restoring an interface
between annuli. A uniform upper bound for the same kernel will make
that estimate stable through any number of annuli.

\subsection{Guard data and interface kernels}

Use the polygon \(P\), the guard, and the end regions of
Section~\ref{sec:annuli}. In lattice units, an annular block has walls
\(RP\) and \(\lambda RP\). We take \(R\) large enough and dyadically
aligned so that all divisions lie between rows of dual vertices.
Every block is connected. The bonds between adjacent blocks are
precisely the bonds to be restored.

Choose one guard vertex in each annulus. Its \emph{guard datum}
\(u\) is the list of all height differences from that vertex to the
other guard vertices, in the chosen connection gauge. Every
compatible datum has positive probability. Denote its law in block
\(i\), in sector \(s\), by \(m_i^s(du)\). Fixing the representative
whose chosen height is zero turns \(u\) into full prescribed height
data on the guard. The configurations on the two sides are then
conditionally independent, since the guard is a full separating
band and the energy is a sum of nearest-neighbor terms.

For two adjacent annuli, with data \(u_i,u_{i+1}\), define
\[
 T_i^s(u_i,u_{i+1})
   =\frac{Z_{\mathrm{restored}}^s(u_i,u_{i+1})}
           {Z_{\mathrm{cut}}^s(u_i,u_{i+1})}.
\]
The numerator is summed modulo one common translation; the
denominator is a product of sums modulo one translation in each
block. Thus the numerator sums the relative integer translation
of the two blocks.

We use the same definition in two other cases. A \emph{core} is
a connected block filling the inner hole of the first annulus;
it may contain a magnetic source and carries no prescribed guard.
Its interface kernel is a function \(T_0^s(u_1)\). At the other
end, several disjoint last annuli may be attached to one connected
exterior, with no direct bonds between these annuli. The exterior
carries no prescribed data. Its one joint kernel is a function of
the incident annular data. In every case let
\[
 d_i=T_i^0(0),
\]
where \(0\) means zero guard differences in the centered sector.
No size estimate for this positive scalar is required.

\begin{lemma}[Upper bound and centered average]\label{kernels:bounds}
For each of the kernels just defined, every sector and every
compatible guard datum satisfy
\begin{equation}
                  0<T_i^s/d_i\le1. \tag{8}\label{eq:source-8}
\end{equation}
For an interface between two annuli, choose the end region \(H\)
on each side containing its seam endpoints. For the joint exterior
kernel choose the outer end region in each incident annulus. Then
\begin{equation}
 1\ge \int (T_i^0/d_i)
                 \prod_{\text{incident annuli }\ell}dm_\ell^0
 \ge \prod_{\text{incident annuli }\ell}
                    e^{-\overline D_\ell(H,p)} .
 \tag{9}\label{eq:source-9}
\end{equation}
\end{lemma}

\begin{proof}
Lemma~\ref{heights:cut-comparison} gives the prescribed-data comparison
\[
 \frac{Z_{\mathrm{restored}}^s(u)}
      {Z_{\mathrm{restored}}^0(0)}
 \le
 \frac{Z_{\mathrm{cut}}^s(u)}{Z_{\mathrm{cut}}^0(0)}.
\]
Rearranging proves Equation~\eqref{eq:source-8}. The comparison
applies to relative guard data: they are affine constraints on
height differences. If auxiliary confinement is used to perform
the comparison, its divergent translation scalar cancels within
each displayed ratio before confinement is removed. The two graphs
need not have the same number of free components.

For the lower bound, work in the centered sector and impose \([H]\)
in every indicated annulus. Adding these centered constraints to
the restored data-versus-zero ratio decreases that ratio. Once they
hold, choose the representative in each annulus to be zero on
its flat end. The restoration factor is then independent of the
guard data: its remaining variables are the relative integer
translations and, in the exterior case, the unprescribed exterior
heights. Let \(\Pbb_{\ell,0}\) denote the free centered law in
annulus \(\ell\). The comparison, after division by the cut
data-versus-zero ratio, gives
\[
 \frac{T_i^0(u)}{d_i}
 \ge
 \prod_{\text{incident annuli }\ell}
 \frac{\Pbb_{\ell,0}([H]\mid u_\ell)}
      {\Pbb_{\ell,0}([H]\mid[p])}.
\]
All seam endpoints lie in the chosen end regions; hence no other
annular heights enter the restoration factor with these events
imposed. Averaging over the independent cut-block laws yields
Equation~\eqref{eq:source-9}.
\end{proof}

\begin{lemma}[Exact multiplication]\label{kernels:multiplication}
Let a core be surrounded by \(N\) annular blocks, with no other
adjacencies between the blocks. Their restored partition function is
\begin{equation}\label{kernels:product}
 Z_{\mathrm{full}}^s
 = Z_{\mathrm{core}}^s\prod_{i=1}^N Z_i^s
   \int T_0^s(u_1)
         \prod_{i=1}^{N-1}T_i^s(u_i,u_{i+1})
         \prod_{i=1}^N dm_i^s(u_i).
\end{equation}
For several disjoint nests attached to a common connected exterior,
the corresponding formula includes its cut partition function and
the one joint exterior kernel.
\end{lemma}

\begin{proof}
Begin with one representative of the heights in each cut block.
On the tree whose vertices are these blocks and whose edges are
their interfaces, integer block constants modulo one common
constant are in bijection with the integer differences along the
tree edges. The bijection has no multiplicity factor: fix the
root constant, then recover each other constant along its unique
path from the root. The relative translations can therefore be
summed independently at the interfaces.

Condition next on every guard datum. Within each annulus, its inner
and outer configurations are conditionally independent. The two
incident interface factors may consequently be integrated separately,
giving the kernels in Equation~\eqref{kernels:product}. For several
nests, put the exterior at the root of the block tree. Its different
boundary portions need not be independent, so their integrations
remain together in the single exterior kernel. This is exactly the
formula asserted.
\end{proof}

\subsection{Small average losses in unit magnetic sectors}

Call an annulus-annulus interface \emph{regular} when the connection
is flat on their restored union and has the same period \(2\pi\),
\(-2\pi\), or zero in both annuli. In particular, restoring that
interface introduces no magnetic source.

\begin{proposition}[Average loss]\label{kernels:average-loss}
For a regular interface,
\begin{equation}
 \int(T_i^s/d_i)\prod_\ell dm_\ell^s
 =
 \left(\int(T_i^0/d_i)\prod_\ell dm_\ell^0\right)
       \frac{F_{\mathrm{restored}}^s}{F_{\mathrm{cut}}^s}.
 \tag{10}\label{eq:source-10}
\end{equation}
Given any \(\xi>0\), first choose \(\varepsilon\) sufficiently small
and fix a polygon \(P\) satisfying its disk inclusions. Then choose
\(\lambda\) sufficiently large. There is a fixed aligned
threshold \(R_0\) such that every regular interface with smaller
inner scale at least \(R_0\) satisfies
\begin{equation}
             \int(1-T_i^s/d_i)\prod_\ell dm_\ell^s\le\xi
 \tag{11}\label{eq:source-11}
\end{equation}
in each of the two unit-period sectors and in the centered sector.
The individual flat-pin interactions in
Equation~\eqref{eq:source-9} can also be made arbitrarily small
with these choices.
\end{proposition}

\begin{proof}
The integral of \(T_i^s\) under the independent cut-block laws is
\(Z_{\mathrm{restored}}^s/Z_{\mathrm{cut}}^s\). Dividing this
identity by its centered counterpart proves
Equation~\eqref{eq:source-10}. The same identity is valid for a
joint exterior kernel, without a claim yet about its average loss.

For a circular annulus of radii \(1\) and \(\lambda\), the minimum
energy with unit period \(2\pi\) is \(2\pi\log\lambda\).
The angular function attains it, and integration on concentric
circles proves the matching lower bound. Restriction of admissible
forms, together with
\(B(0,1-\varepsilon)\subset P\subset B(0,1+\varepsilon)\), gives
\[
 \mathcal M(\lambda P\setminus P^\circ)
       =2\pi\log\lambda+O(\varepsilon),
\]
where the constant is independent of large \(\lambda\). For example,
the circular annulus with radii \(1+\varepsilon\) and
\(\lambda(1-\varepsilon)\) is contained in the polygonal one,
which in turn is contained in the circular annulus with radii
\(1-\varepsilon\) and \(\lambda(1+\varepsilon)\).
The same bounds apply with \(\lambda^2\) to the union of two blocks.
It follows that
\[
 0\le
 \mathcal M(\lambda^2P\setminus P^\circ)
       -2\mathcal M(\lambda P\setminus P^\circ)
 \le C\varepsilon.
\]
The lower bound follows by restriction to the two cut annuli;
their energies agree by scale invariance in dimension two.
Theorem~\ref{sectors:positive-radius}, applied at this fixed geometry,
therefore makes the last factor in Equation~\eqref{eq:source-10}
converge to a number in \([e^{-C\varepsilon/(2a)},1]\).
Changing the unit period's sign does not change its energy.

Choose \(\varepsilon\) to make this error small, and then choose
\(\lambda\) so that Proposition~\ref{annuli:flat-interaction} makes
the two centered pin interactions small.
Equation~\eqref{eq:source-9} and
Equation~\eqref{eq:source-10} now give the desired small loss in
the mesh limit. With \(P\) and \(\lambda\) fixed, every regular
interface is a refinement of one fixed geometry. The convergence
theorems hold along every admissible mesh sequence, so there is
one threshold \(R_0\) beyond which the inequalities hold at every
aligned scale. There are only three sectors to consider.
This proves Equation~\eqref{eq:source-11}.
\end{proof}

\subsection{A positive-kernel stability lemma}

The preceding estimate controls the average loss under independent
guard laws. Actual guard data in a nest have been influenced by
all the inner blocks. The following lemma bounds that change of law
without accumulating a fixed error at every interface.

\begin{lemma}[Stability from a positive initial weight]
\label{kernels:stability}
Let \((X_i,m_i)\), \(i\ge1\), be probability spaces, and let
\(t_i:X_i\times X_{i+1}\to(0,1]\) be measurable kernels with
\[
 \int(1-t_i)\,d(m_i\otimes m_{i+1})\le\xi,\qquad
                         0<\xi<1/64.
\]
Fix a measurable \(g:X_1\to(0,1]\), positive \(m_1\)-almost
everywhere. For each \(N\), give \(X_1\times\cdots\times X_N\)
density proportional to
\[
                    g(u_1)\prod_{i=1}^{N-1}t_i(u_i,u_{i+1})
\]
with respect to \(\prod_i m_i\). For all sufficiently large \(N\),
its last-coordinate marginal has density at most \(2\) with respect
to \(m_N\). No positive essential lower bound on \(g\) is required.
\end{lemma}

\begin{proof}
Define
\[
 L_i f(v)=\int t_i(u,v)f(u)\,dm_i(u).
\]
First propagate the constant density:
\[
 q_1=1,\qquad
 \beta_i=\int L_iq_i\,dm_{i+1},\qquad
 q_{i+1}=\beta_i^{-1}L_iq_i.
\]
Inductively \(q_i\le2\), because
\[
 \beta_i
  =1-\int(1-t_i)q_i\,d(m_i\otimes m_{i+1})
  \ge1-2\xi,\qquad
 q_{i+1}\le M:=(1-2\xi)^{-1}<2.
\]
The same argument, with the kernels reversed, bounds by \(2\) any
normalized density obtained by propagating \(1\) backwards through
a finite number of them.

Propagate \(g\) using these same scalars, rather than renormalizing
it at each step:
\[
 f_1=g,\qquad f_{i+1}=\beta_i^{-1}L_if_i,\qquad
 c_i=\int f_i\,dm_i,\qquad E_i=f_i-c_iq_i.
\]
The scalar \(c_i\) is the ratio of the two integrals obtained by
starting with \(g\) and with \(1\). Equivalently, it is the
expectation of \(g\) under a normalized backwards density on
\(X_1\), which is at most \(2\). Choose \(\delta>0\) with
\(m_1(g<\delta)<1/4\). Every such density assigns at least half
its mass to \(\{g\ge\delta\}\), so
\begin{equation}\label{kernels:core-mass}
                         c_i\ge\delta/2
                    \quad\text{for every }i.
\end{equation}

For a mean-zero function \(E\), replace \(t_i\) in \(L_iE\) by
\(t_i-1\). Since \(0\le(1-t_i)^2\le1-t_i\), the resulting
operator has \(L^2(m_i)\)-to-\(L^2(m_{i+1})\) norm at most
\(\sqrt\xi\). The exact recurrence for the centered remainder is
\[
 E_{i+1}
 =\frac{L_iE_i-\bigl(\int L_iE_i\,dm_{i+1}\bigr)q_{i+1}}
        {\beta_i}.
\]
Using \(\|q_{i+1}\|_2\le2\), we obtain
\begin{equation}\label{kernels:contraction}
 \|E_{i+1}\|_2
 \le\frac{3\sqrt\xi}{1-2\xi}\|E_i\|_2.
\end{equation}
The factor is strictly smaller than one. Also \(t_i\le1\) gives
\(\|L_iE_i\|_\infty\le\|E_i\|_2\), and hence
\[
 \|E_{i+1}\|_\infty
       \le\frac{3}{1-2\xi}\|E_i\|_2.
\]
Thus the remainders tend to zero in the last displayed norm.
The last-coordinate density in the lemma is
\[
                    f_N/c_N=q_N+E_N/c_N.
\]
Equation~\eqref{kernels:core-mass} and the strict inequality
\(M<2\) show that it is eventually bounded by \(2\).
\end{proof}

\subsection{Application to a nest of annuli}

Fix the parameters in Proposition~\ref{kernels:average-loss} with
\(\xi<1/64\), and increase the aligned \(R_0\) if necessary. Set
\[
                 R_j=R_0\lambda^j,\qquad j\ge0.
\]
Let annulus \(i\) have walls \(R_{i-1}P\) and \(R_iP\), and fill
the inner hole by the fixed core in \(R_0P\). The core is a fixed
finite graph in lattice units. In a unit magnetic sector it
contains the source; all regular annuli carry its period.
Fix the core and the connection on the increasing family of blocks
once and for all, so restrictions agree as the depth \(N\) increases.
A fixed straight connection cut from the source gives such a choice.
Thus the first guard law and its core interface kernel do not vary
with \(N\).

\begin{corollary}[The last guard law]\label{kernels:last-guard}
For these fixed choices, the last-guard law in a nest of \(N\)
annuli has density at most \(2\) relative to \(m_N^s\) for all
sufficiently large \(N\), in the centered sector and in either
unit-period sector. The conclusion holds for identical lattice
translates of the nest with the same threshold in \(N\).
\end{corollary}

\begin{proof}
Put \(t_i=T_i^s/d_i\) for regular interfaces and
\(g=T_0^s/d_0\) for the core interface. Lemma~\ref{kernels:bounds}
gives \(0<g\le1\) and \(0<t_i\le1\).
Proposition~\ref{kernels:average-loss} gives the average-loss
hypothesis of Lemma~\ref{kernels:stability}.
Exact multiplication in Lemma~\ref{kernels:multiplication}
identifies its last marginal with the actual last-guard law.
The fixed initial weight \(g\) may be arbitrarily small on some
data; positivity is sufficient.

A lattice translation preserves all block laws and kernels.
Changing the source sign negates the connection and heights.
More generally, a connection with the same source restricts on
the nest to the chosen connection up to an integral gauge change,
which merely reindexes the guard data. Taking the maximum of the
finitely many thresholds for the centered and two unit sectors
proves the assertion.
\end{proof}
